\section{Идея SVM: зачем он нужен и почему не хватает логистической регрессии}

После детального рассмотрения линейно-разделимых и неразделимых случаев естественным образом возникает фундаментальный вопрос о месте метода опорных векторов в современном арсенале алгоритмов машинного обучения. Действительно, когда существует такая классическая и хорошо изученная модель как логистическая регрессия, зачем разрабатывать ещё один метод для решения задачи бинарной классификации?

\subsection{Различие двух подходов}

Если внимательно проанализировать теоретические основы логистической регрессии и метода опорных векторов, становится ясно, что хотя оба метода решают задачу классификации, они исходят из совершенно различных предпосылок и оптимизируют разные целевые функции. 

Логистическая регрессия нацелена на максимизацию правдоподобия наблюдаемых данных. Её математический аппарат строится вокруг моделирования вероятности принадлежности объекта к положительному классу с помощью сигмоидной функции. Формально это выражается как $P(y = +1 | x) = \sigma(w^\top x + b)$. (вероятность положительного класса на текущем объекте $x$). Процесс обучения в этом случае сводится к максимизации логарифмического правдоподобия или, что эквивалентно, минимизации логистической функции потерь.

Метод опорных векторов, напротив, реализует принципиально иную философию. Вместо того чтобы стремиться к идеальной классификации всех точек, SVM фокусируется исключительно на тех наблюдениях, которые находятся ближе всего к границе между классами — так называемых опорных векторах. Основная цель алгоритма — не минимизация ошибки на всей выборке, а максимизация зазора между классами, то есть расстояния от разделяющей гиперплоскости до ближайших точек каждого класса.

\subsection{Почему максимизация зазора важна?}

Максимизация зазора между классами — это не просто математическая элегантность или эстетическое предпочтение. Данный принцип имеет глубокое теоретическое обоснование и важные практические следствия, связанные с обобщающей способностью модели.

Логистическая регрессия, стремясь правильно классифицировать все точки без исключения, может существенно изменить положение разделяющей прямой, чтобы учесть выбросы. В результате граница решения окажется расположенной близко к основной массе точек обоих классов, что потенциально ухудшит обобщающую способность модели на новых данных. В этом случае SVM показывает себя значительно устойчивее на выбросах, т.к. ориентируется в основном на опроные векторы.

И это всё приводит к занятной мысли "--- обычно, у SVM базовые метрики ($F$"=мера, $precision$, $recall$...) выше, чем у логистической регрессии. С другой стороны, логистическая регрессия позволяет корректно работать с вероятностями. Модель подбирается по нужде, и мы ещё успеем убедиться, что не существует универсального алгоритма, который способен решать задачу лучше всех остальных.
